\begin{problem}[Finite principal-cover criterion]\label{prob:vi18-p06}
For $f_1,\dots,f_r\in A$, prove $\Spec A=\bigcup_iD(f_i)$ iff $(f_1,\dots,f_r)=A$, and formulate the infinite-family version.
\end{problem}

\ifdefined\FullProblemDossiers
\paragraph{Why this problem matters.}
This problem consolidates the cover criterion aspect of affine geometry and turns a structural statement into a reusable calculation.

\paragraph{Useful ideas before starting.}
Use complements and maximal ideals.

\paragraph{Guided hints.}
\begin{enumerate}
\item Use complements and maximal ideals.
\item Reduce the statement to a ring-theoretic property whenever possible.
\item Translate the result back into topology, sheaves, or local rings so that all affine-scheme layers are visible.
\end{enumerate}

\begin{solution}
Put $I=(f_1,\dots,f_r)$. The complement of the union of distinguished opens is
\[
\Spec A\setminus\bigcup_{i=1}^rD(f_i)
=\bigcap_{i=1}^rV(f_i)=V(I).
\]
Thus the $D(f_i)$ cover $\Spec A$ iff $V(I)=\varnothing$.

If $I=A$, then $V(I)=V(1)=\varnothing$. Conversely, if $I$ is proper, the maximal-ideal theorem gives a maximal ideal $\mathfrak m\supseteq I$. Then $\mathfrak m\in V(I)$, so $V(I)$ is nonempty. Hence
\[
\Spec A=\bigcup_{i=1}^rD(f_i)
\iff I=A
\iff 1\in(f_1,\dots,f_r).
\]
Equivalently there exist $a_i\in A$ with $1=\sum_i a_if_i$.

For an arbitrary family $\{f_i\}_{i\in I}$, the same proof gives
\[
\Spec A=\bigcup_{i\in I}D(f_i)
\iff (f_i\mid i\in I)=A.
\]
Because membership of $1$ in the generated ideal uses only finitely many generators, an infinite distinguished-open cover automatically contains a finite distinguished-open subcover.
\end{solution}


\paragraph{What happened mathematically.}
The calculation illustrates that global affine geometry is controlled by commutative algebra once the structure sheaf is kept in view.

\paragraph{Extensions.}
Vary the ring by products, quotients, localizations, arithmetic examples, or nilpotent thickenings and repeat the same dictionary.

\paragraph{Provenance.}
\textbf{New canonical synthesis; no numbered legacy Problem is claimed.}
\else
\paragraph{Provenance.} \textbf{New canonical synthesis; no numbered legacy Problem is claimed.}
\fi
